\documentclass[10pt,a4paper]{amsart}
\usepackage{amsmath,amssymb,amsthm,graphicx}
% Retained minimal preamble for the Stacks Project chapter flat.tex (More on Flatness).
% Verbatim excerpt of the chapter preamble: Xy-pic setup (native preamble lines 14-18) and the native
% theorem declarations with the mathematical macros (native preamble lines 154-217). Copied mechanically;
% nothing added, removed or redefined. The wrapper class and the amsmath/amssymb/amsthm/graphicx base
% packages are supplied by the assembler, so the class/externaldocument part of the native preamble is omitted.
% For commutative diagrams we use Xy-pic
\usepackage[all]{xy}

% We use 2cell for 2-commutative diagrams.
\xyoption{2cell}
% Theorem environments.
%
\theoremstyle{plain}
\newtheorem{theorem}[subsection]{Theorem}
\newtheorem{proposition}[subsection]{Proposition}
\newtheorem{lemma}[subsection]{Lemma}

\theoremstyle{definition}
\newtheorem{definition}[subsection]{Definition}
\newtheorem{example}[subsection]{Example}
\newtheorem{exercise}[subsection]{Exercise}
\newtheorem{situation}[subsection]{Situation}

\theoremstyle{remark}
\newtheorem{remark}[subsection]{Remark}
\newtheorem{remarks}[subsection]{Remarks}

\numberwithin{equation}{subsection}

% Macros
%
\def\lim{\mathop{\mathrm{lim}}\nolimits}
\def\colim{\mathop{\mathrm{colim}}\nolimits}
\def\Spec{\mathop{\mathrm{Spec}}}
\def\Hom{\mathop{\mathrm{Hom}}\nolimits}
\def\Ext{\mathop{\mathrm{Ext}}\nolimits}
\def\SheafHom{\mathop{\mathcal{H}\!\mathit{om}}\nolimits}
\def\SheafExt{\mathop{\mathcal{E}\!\mathit{xt}}\nolimits}
\def\Sch{\mathit{Sch}}
\def\Mor{\mathop{\mathrm{Mor}}\nolimits}
\def\Ob{\mathop{\mathrm{Ob}}\nolimits}
\def\Sh{\mathop{\mathit{Sh}}\nolimits}
\def\NL{\mathop{N\!L}\nolimits}
\def\CH{\mathop{\mathrm{CH}}\nolimits}
\def\proetale{{pro\text{-}\acute{e}tale}}
\def\etale{{\acute{e}tale}}
\def\QCoh{\mathit{QCoh}}
\def\Ker{\mathop{\mathrm{Ker}}}
\def\Im{\mathop{\mathrm{Im}}}
\def\Coker{\mathop{\mathrm{Coker}}}
\def\Coim{\mathop{\mathrm{Coim}}}

% Boxtimes
%
\DeclareMathSymbol{\boxtimes}{\mathbin}{AMSa}{"02}

%
% Macros for moduli stacks/spaces
%
\def\QCohstack{\mathcal{QC}\!\mathit{oh}}
\def\Cohstack{\mathcal{C}\!\mathit{oh}}
\def\Spacesstack{\mathcal{S}\!\mathit{paces}}
\def\Quotfunctor{\mathrm{Quot}}
\def\Hilbfunctor{\mathrm{Hilb}}
\def\Curvesstack{\mathcal{C}\!\mathit{urves}}
\def\Polarizedstack{\mathcal{P}\!\mathit{olarized}}
\def\Complexesstack{\mathcal{C}\!\mathit{omplexes}}
% \Pic is the operator that assigns to X its picard group, usage \Pic(X)
% \Picardstack_{X/B} denotes the Picard stack of X over B
% \Picardfunctor_{X/B} denotes the Picard functor of X over B
\def\Pic{\mathop{\mathrm{Pic}}\nolimits}
\def\Picardstack{\mathcal{P}\!\mathit{ic}}
\def\Picardfunctor{\mathrm{Pic}}
\def\Deformationcategory{\mathcal{D}\!\mathit{ef}}
\usepackage{cleveref}
\crefname{originallemma}{original lemma}{original lemmas}
\Crefname{originallemma}{Original lemma}{Original lemmas}
\makeatletter
\newlabel{lemma-elementary-devissage-variant}{{05H6}{source}{}{sourcepointer-1}{}}
\newlabel{lemma-elementary-devissage-variant@cref}{{[originallemma][0][]05H6}{[1][0][]source}{}{}{}}
\makeatother
\crefname{originallemma}{original lemma}{original lemmas}
\Crefname{originallemma}{Original lemma}{Original lemmas}
\makeatletter
\newlabel{lemma-existence-alpha}{{05HF}{source}{}{sourcepointer-2}{}}
\newlabel{lemma-existence-alpha@cref}{{[originallemma][0][]05HF}{[1][0][]source}{}{}{}}
\makeatother
\crefname{originallemma}{original lemma}{original lemmas}
\Crefname{originallemma}{Original lemma}{Original lemmas}
\makeatletter
\newlabel{lemma-shrink}{{05HA}{source}{}{sourcepointer-3}{}}
\newlabel{lemma-shrink@cref}{{[originallemma][0][]05HA}{[1][0][]source}{}{}{}}
\makeatother
\crefname{originalreference}{original reference}{original references}
\Crefname{originalreference}{Original reference}{Original references}
\makeatletter
\newlabel{item-shrink-base}{{05HA(1)}{source}{}{sourcepointer-4}{}}
\newlabel{item-shrink-base@cref}{{[originalreference][0][]05HA(1)}{[1][0][]source}{}{}{}}
\makeatother
\crefname{originalreference}{original reference}{original references}
\Crefname{originalreference}{Original reference}{Original references}
\makeatletter
\newlabel{item-shrink-on-Y}{{05HA(2)}{source}{}{sourcepointer-5}{}}
\newlabel{item-shrink-on-Y@cref}{{[originalreference][0][]05HA(2)}{[1][0][]source}{}{}{}}
\makeatother
\crefname{originallemma}{original lemma}{original lemmas}
\Crefname{originallemma}{Original lemma}{Original lemmas}
\makeatletter
\newlabel{lemma-elementary-etale-neighbourhood}{{05HE}{source}{}{sourcepointer-6}{}}
\newlabel{lemma-elementary-etale-neighbourhood@cref}{{[originallemma][0][]05HE}{[1][0][]source}{}{}{}}
\makeatother
\crefname{originallemma}{original lemma}{original lemmas}
\Crefname{originallemma}{Original lemma}{Original lemmas}
\makeatletter
\newlabel{lemma-base-change-complete-at-x}{{05HK}{source}{}{sourcepointer-7}{}}
\newlabel{lemma-base-change-complete-at-x@cref}{{[originallemma][0][]05HK}{[1][0][]source}{}{}{}}
\makeatother
\crefname{originallemma}{original lemma}{original lemmas}
\Crefname{originallemma}{Original lemma}{Original lemmas}
\makeatletter
\newlabel{lemma-shrink-complete}{{05HM}{source}{}{sourcepointer-8}{}}
\newlabel{lemma-shrink-complete@cref}{{[originallemma][0][]05HM}{[1][0][]source}{}{}{}}
\makeatother
\crefname{originalreference}{original reference}{original references}
\Crefname{originalreference}{Original reference}{Original references}
\makeatletter
\newlabel{item-shrink-on-X-complete}{{05HM(3)}{source}{}{sourcepointer-9}{}}
\newlabel{item-shrink-on-X-complete@cref}{{[originalreference][0][]05HM(3)}{[1][0][]source}{}{}{}}
\makeatother
\title{Existence of a complete devissage at a point, by induction on the dimension of the fibre support}
\author{The Stacks Project Authors}
\date{stacks-project:flat.tex:local-snapshot-20261011; GNU FDL 1.2 or later}
\begin{document}
\maketitle
\noindent\textit{Editorial note.} All delivered mathematical text is the author's own, copied from the pinned Stacks Project chapter file flat.tex; no statement or proof step was rewritten, reordered or removed. The statement writes 'has a complete devissage at x' while the proof works with x'; both are reproduced verbatim.\par\medskip
\noindent\textit{Editorial note.} The preamble is the chapter's own retained minimal preamble (sources/stacks-flat/retained-preamble.tex), a verbatim excerpt of the native Xy-pic setup and of the native theorem declarations and mathematical macros; the Xy-pic xy environment is required because the counted statement and proof display four commutative diagrams. No macro was invented or redefined.\par\medskip
\noindent\textit{Editorial note.} The two displayed equations of the proof (equation-overcome-this and equation-shrink-this) are labelled and referenced inside the counted spans, so they are resolved by the document itself and are not listed as external references.\par\medskip
\noindent\textit{Editorial note.} Same-chapter referenced results are located at their \textbackslash{}label line inside the pinned chapter file; the four enumerate items are located at their own \textbackslash{}label lines and are recorded with the tag of the lemma that owns them together with the item index, because the Stacks tag list assigns tags to results and not to the individual items of an enumerate.\par\medskip
\noindent\textit{Allowed standard prerequisites (editorial).}\begin{itemize}
\item Complete devissage of a finite type module. Same chapter, Definition environment labelled definition-complete-devissage, source line 1195 (Stacks tag 05HH). Supplies what the proposition has to produce, i.e. the tuple (Z\_k, Y\_k, i\_k, pi\_k, G\_k, alpha\_k) with its conditions over the base point; the induction itself is the counted content.
\item One step devissage and the standard shrinking operations. Same chapter, Definition definition-one-step-devissage (05H4, line 780), Definition definition-shrink (05H9, line 850) and the items of Lemma lemma-shrink (05HA, line 970) and of Lemma lemma-shrink-complete (05HM, line 1361). Supply the vocabulary of standard shrin kings and the three listed shrinking moves (shrink the base, shrink so that Y' = W x\_S S', shrink so that X' is inside a given open) that the counted proof invokes by item number; the way they are iterated to make the two devissages fit together stays inside the counted span.
\item Pointwise one step devissage, its finite-presentation and flatness behaviour, and the affine base change of complete devissages. Same chapter, Lemma lemma-elementary-devissage-variant (05H6, line 855), Lemma lemma-existence-alpha (05HF, line 1138), Lemma lemma-elementary-etale-neighbourhood (05HE, line 1027) and Lemma lemma-base-change-complete-at-x (05HK, line 1313). Supply the four inputs of the induction step: the one step devissage after elementary etale localisation, the free cover alpha\_1 with a strictly smaller cokernel on the fibre, the diagram that transports a one step devissage up to a complete one, and the affine base change of the existing complete devissage.
\end{itemize}\medskip
\ifdefined\hypertarget\hypertarget{sourcepointer-1}{}\fi
\noindent\textit{Source reference.} Original lemma 05H6: sources/stacks-flat/pinned-original-source.tex, source lines 855--855. Same-chapter pointer invoked at proof line 1450 to reduce to the case where a one step devissage exists; the induction, the dimension descent and the final gluing are retained verbatim in the delivered body.\par\medskip
\ifdefined\hypertarget\hypertarget{sourcepointer-2}{}\fi
\noindent\textit{Source reference.} Original lemma 05HF: sources/stacks-flat/pinned-original-source.tex, source lines 1138--1138. Same-chapter pointer invoked at proof line 1468 to produce alpha\_1 with the strict dimension drop; it supplies the input of the induction step, while the case split on the stalk of the cokernel and the induction itself are retained in the counted span.\par\medskip
\ifdefined\hypertarget\hypertarget{sourcepointer-3}{}\fi
\noindent\textit{Source reference.} Original lemma 05HA: sources/stacks-flat/pinned-original-source.tex, source lines 970--970. Same-chapter pointer to the shrinking lemma, invoked at proof lines 1482, 1552, 1560 and 1562; the iteration of its items that makes the one step and complete devissages fit together is the counted concluding step of this proof.\par\medskip
\ifdefined\hypertarget\hypertarget{sourcepointer-4}{}\fi
\noindent\textit{Source reference.} Original reference 05HA(1): sources/stacks-flat/pinned-original-source.tex, source lines 976--976. Same-chapter item pointer invoked at proof line 1560 for the move 'shrink the base to S''; the item itself is not reproduced, and the way it is used is retained in the counted span.\par\medskip
\ifdefined\hypertarget\hypertarget{sourcepointer-5}{}\fi
\noindent\textit{Source reference.} Original reference 05HA(2): sources/stacks-flat/pinned-original-source.tex, source lines 981--981. Same-chapter item pointer invoked at proof lines 1482 and 1562 for the move that realises Y''\_1 as W x\_S S''; its use in arranging Y''\_1 = X'\_1 x\_S S'' is retained in the counted span.\par\medskip
\ifdefined\hypertarget\hypertarget{sourcepointer-6}{}\fi
\noindent\textit{Source reference.} Original lemma 05HE: sources/stacks-flat/pinned-original-source.tex, source lines 1027--1027. Same-chapter pointer invoked at proof line 1508 to build the diagram in which the one step devissage and the complete devissage of the cokernel are assembled; the subsequent base change and shrinking are retained in the counted span.\par\medskip
\ifdefined\hypertarget\hypertarget{sourcepointer-7}{}\fi
\noindent\textit{Source reference.} Original lemma 05HK: sources/stacks-flat/pinned-original-source.tex, source lines 1313--1313. Same-chapter pointer invoked at proof line 1524 to transport the complete devissage of the cokernel to S''; the resulting reorganisation of the data is retained in the counted span.\par\medskip
\ifdefined\hypertarget\hypertarget{sourcepointer-8}{}\fi
\noindent\textit{Source reference.} Original lemma 05HM: sources/stacks-flat/pinned-original-source.tex, source lines 1361--1361. Same-chapter pointer to the shrinking lemma for complete devissages, invoked at proof lines 1552 and 1555; the two-stage shrinking that concludes the proof is retained in the counted span.\par\medskip
\ifdefined\hypertarget\hypertarget{sourcepointer-9}{}\fi
\noindent\textit{Source reference.} Original reference 05HM(3): sources/stacks-flat/pinned-original-source.tex, source lines 1376--1376. Same-chapter item pointer invoked at proof line 1555 to shrink the complete devissage so that X'\_1 lies inside Y\_1; the item itself is not reproduced, and the use made of it is retained in the counted span.\par\medskip
\begin{proposition}
\label{proposition-existence-complete-at-x}
Let $S$ be a scheme.
Let $X$ be locally of finite type over $S$.
Let $x \in X$ be a point with image $s \in S$.
There exists a commutative diagram
$$
\xymatrix{
(X, x) \ar[d] & (X', x') \ar[l]^g \ar[d] \\
(S, s) & (S', s') \ar[l]
}
$$
of pointed schemes such that the horizontal
arrows are elementary \'etale neighbourhoods
and such that $g^*\mathcal{F}/X'/S'$ has a complete
d\'evissage at $x$.
\end{proposition}

\begin{proof}
We prove this by induction on the integer
$d = \dim_x(\text{Supp}(\mathcal{F}_s))$.
By
Lemma \ref{lemma-elementary-devissage-variant}
there exists a diagram
$$
\xymatrix{
(X, x) \ar[d] & (X', x') \ar[l]^g \ar[d] \\
(S, s) & (S', s') \ar[l]
}
$$
of pointed schemes such that the horizontal
arrows are elementary \'etale neighbourhoods
and such that $g^*\mathcal{F}/X'/S'$ has a one step d\'evissage at $x'$.
The local nature of the problem implies that we may replace
$(X, x) \to (S, s)$ by $(X', x') \to (S', s')$. Thus after doing so
we may assume that there exists a one step d\'evissage
$(Z_1, Y_1, i_1, \pi_1, \mathcal{G}_1)$ of $\mathcal{F}/X/S$ at $x$.

\medskip\noindent
We apply
Lemma \ref{lemma-existence-alpha}
to find a map
$$
\alpha_1 :
\mathcal{O}_{Y_1}^{\oplus r_1}
\longrightarrow
\pi_{1, *}\mathcal{G}_1
$$
which induces an isomorphism of vector spaces over $\kappa(\xi_1)$
where $\xi_1 \in Y_1$ is the unique generic point of the fibre of
$Y_1$ over $s$. Moreover
$\dim_{y_1}(\text{Supp}(\Coker(\alpha_1)_s)) < d$.
It may happen that the stalk of $\Coker(\alpha_1)_s$
at $y_1$ is zero. In this case we may shrink $Y_1$ by
Lemma \ref{lemma-shrink} (\ref{item-shrink-on-Y})
and assume that $\Coker(\alpha_1) = 0$ so we obtain a
complete d\'evissage of length zero.

\medskip\noindent
Assume now that the stalk of $\Coker(\alpha_1)_s$
at $y_1$ is not zero. In this case, by induction, there exists a
commutative diagram
\begin{equation}
\label{equation-overcome-this}
\vcenter{
\xymatrix{
(Y_1, y_1) \ar[d] & (Y'_1, y'_1) \ar[l]^h \ar[d] \\
(S, s) & (S', s') \ar[l]
}
}
\end{equation}
of pointed schemes such that the horizontal
arrows are elementary \'etale neighbourhoods
and such that $h^*\Coker(\alpha_1)/Y'_1/S'$ has a complete
d\'evissage
$$
(Z_k, Y_k, i_k, \pi_k, \mathcal{G}_k, \alpha_k, z_k, y_k)_{k = 2, \ldots, n}
$$
at $y'_1$. (In particular $i_2 : Z_2 \to Y'_1$ is a closed immersion into
$Y'_2$.) At this point we apply
Lemma \ref{lemma-elementary-etale-neighbourhood}
to $S, X, \mathcal{F}, x, s$, the system
$(Z_1, Y_1, i_1, \pi_1, \mathcal{G}_1)$ and
diagram (\ref{equation-overcome-this}). We obtain a diagram
$$
\xymatrix{
& & (X'', x'') \ar[lld] \ar[d] & (Z''_1, z''_1) \ar[l] \ar[lld] \ar[d] \\
(X, x) \ar[d] & (Z_1, z_1) \ar[l] \ar[d] &
(S'', s'') \ar[lld] & (Y''_1, y''_1) \ar[lld] \ar[l] \\
(S, s) & (Y_1, y_1) \ar[l]
}
$$
with all the properties as listed in the referenced lemma.
In particular $Y''_1 \subset Y'_1 \times_{S'} S''$. Set
$X_1 = Y'_1 \times_{S'} S''$ and let $\mathcal{F}_1$ denote the
pullback of $\Coker(\alpha_1)$. By
Lemma \ref{lemma-base-change-complete-at-x}
the system
\begin{equation}
\label{equation-shrink-this}
(Z_k \times_{S'} S'',
Y_k \times_{S'} S'', i''_k, \pi''_k, \mathcal{G}''_k,
\alpha''_k, z''_k, y''_k)_{k = 2, \ldots, n}
\end{equation}
is a complete d\'evissage of $\mathcal{F}_1$
to $X_1$. Again, the nature of the problem allows
us to replace $(X, x) \to (S, s)$ by $(X'', x'') \to (S'', s'')$.
In this we see that we may assume:
\begin{enumerate}
\item[(a)] There exists a one step d\'evissage
$(Z_1, Y_1, i_1, \pi_1, \mathcal{G}_1)$ of $\mathcal{F}/X/S$ at $x$,
\item[(b)] there exists an $\alpha_1 : \mathcal{O}_{Y_1}^{\oplus r_1}
\to \pi_{1, *}\mathcal{G}_1$ such that $\alpha \otimes \kappa(\xi_1)$
is an isomorphism,
\item[(c)] $Y_1 \subset X_1$ is open, $y_1 = x_1$, and
$\mathcal{F}_1|_{Y_1} \cong \Coker(\alpha_1)$, and
\item[(d)] there exists a complete d\'evissage
$(Z_k, Y_k, i_k, \pi_k, \mathcal{G}_k, \alpha_k, z_k, y_k)_{k = 2, \ldots, n}$
of $\mathcal{F}_1/X_1/S$ at $x_1$.
\end{enumerate}
To finish the proof all we have to do is shrink the one step d\'evissage
and the complete d\'evissage such that they fit together to a complete
d\'evissage. (We suggest the reader do this on their own using
Lemmas \ref{lemma-shrink} and
\ref{lemma-shrink-complete}
instead of reading the proof that follows.) Since $Y_1 \subset X_1$
is an open neighbourhood of $x_1$ we may apply
Lemma \ref{lemma-shrink-complete} (\ref{item-shrink-on-X-complete})
to find a standard shrinking $S', X'_1, Z'_2, Y'_2, \ldots, Y'_n$
of the datum (d) so that $X'_1 \subset Y_1$. Note that $X'_1$ is also
a standard open of the affine scheme $Y_1$. Next, we shrink the datum
(a) as follows: first we shrink the base $S$ to $S'$, see
Lemma \ref{lemma-shrink} (\ref{item-shrink-base}) and then
we shrink the result to $S''$, $X''$, $Z''_1$, $Y''_1$ using
Lemma \ref{lemma-shrink} (\ref{item-shrink-on-Y})
such that eventually $Y''_1 = X'_1 \times_S S''$ and $S'' \subset S'$.
Then we see that
$$
Z''_1, Y''_1, Z'_2 \times_{S'} S'', Y'_2 \times_{S'} S'', \ldots,
Y'_n \times_{S'} S''
$$
gives the complete d\'evissage we were looking for.
\end{proof}

\end{document}
